
%%%%%%%%%%%%%%%%%%%%%%%%%%%%%%%%%%%%
%%
%%
%%%%%%%%%%%%%%%%%%%%%%%%%%%%%%%%%%%%
\section{Useful results about pseudo-Lipschitz functions and CGMT}\label{SM useful fact}
For $k\geq 1$ we say a function $f:\R^m\rightarrow\R$ is pseudo-Lipschitz of order $k$ and denote it by $f\in \rm{PL}(k)$ if there exists a constant $L>0$ such that, for all $\x,\y\in\R^m$:
\begin{align}
|f(\x)-f(\y)|\leq L\left(1+\|\x\|_2^{k-1}+\|\y\|_2^{k-1}\right)\|\x-\y\|_2.\label{PL func}
\end{align}
In particular, when $f\in\rm{PL}(k)$, the following properties hold; see \cite{bayati2011dynamics}:
\begin{enumerate}
\item There exists a constant $L'$ such that for all $\x\in\R^n$: $|f(\x)|\leq L'(1+\|\x\|_2^k).$

\item $f$ is locally Lipschitz, that is for any $M>0$, there exists a constant $L_{M,m}<\infty$ such that for all $x,y\in[-M,M]^m$,
$
|f(\x)-f(\y)| \leq L_{M,m}\|\x-\y\|_2.
$
Further, $L_{M,m}\leq c(1+(M\sqrt{m})^{k-1})$ for some costant $c$.
\end{enumerate}

Using the above properties, we prove the following two technical lemmas used in the proof of Theorem \ref{thm:master_W2}

\begin{lemma}\label{lem:fL}
Let $g:\R\rightarrow\R$ be a Lipschitz function. Consider the function $f:\R^3\rightarrow\R$ defined as follows:
$$
f(\x) = x_3(x_2-g(x_1))^2.
$$
Then, $f\in\rm{PL}(3)$. If additionally, $f:\R^2\times \Zc\rightarrow\R$ for a bounded set $\Zc\subset\R$, then $f\in\Plt\subset\rm{PL}(3)$. Specifically, setting $\Zc=[\Sigma_{\min},\Sigma_{\max}]$ (as per Assumption \ref{ass:mu}), we find that $\Fl\subset\Plt$, where $\Fl$ is defined in \eqref{eq:fdef}.
\end{lemma}

\begin{proof}
We first prove that $f\in\rm{PL}(3)$.

Let $h:\R^2\rightarrow\R$ defined as $h(\ub)=(\ub_2-g(\ub_1))^2$. The function $(\ub_1,\ub_2)\mapsto\ub_2-g(\ub_1)$ is clearly Lipschitz. Thus, $h\in\rm{PL}(2)$, i.e., 
\begin{align}
|h(\ub)-h(\vb)| \leq C(1+\|\ub\|_2+\|\vb\|_2)\|\ub-\vb\|_2\quad\text{and}\quad |h(\vb)|\leq C'(1+\|\vb\|_2^2).\label{eq:h_pl}
\end{align} 
Therefore, letting $\x=(\ub,x_3)\in\R^3$ and $\y=(\vb,y_3)\in\R^3$, we have that 
\begin{align}
|f(\x)-f(\y)| &= |x_3h(\ub) - y_3h(\vb)| \leq |x_3||h(\ub)-h(\vb)| + |h(\vb)| |x_3-y_3|\nn\\
%&\leq |x_3|(1+\|\ub\|_2+\|\vb\|_2)\|\ub-\vb\|_2 + |h(\vb)| |x_3-y_3| \nn\\
&\leq C|x_3|(1+\|\ub\|_2+\|\vb\|_2)\|\ub-\vb\|_2 + C'(1+\|\vb\|_2^2)|x_3-y_3| \nn\\
&\leq C(|x_3|^2+(1+\|\ub\|_2+\|\vb\|_2)^2)\|\ub-\vb\|_2 + C'(1+\|\vb\|_2^2)|x_3-y_3| \nn\\
&\leq C(1+\|\x\|_2^2+\|\y\|_2^2)\|\ub-\vb\|_2 + C'(1+\|\x\|_2^2+\|\y\|_2^2)|x_3-y_3| \nn\\
&\leq C(1+\|\x\|_2^2+\|\y\|_2^2)\|\x-\y\|_2.
\end{align}
In the second line, we used \eqref{eq:h_pl}. In the third line, we used $2xy\leq x^2+y^2$. In the fourth line, we used Cauchy-Schwarz inequality. $C,C'>0$ are absolute constants that may change from line to line.

\fy{Now, we shall prove that $f\in\Plt$. To accomplish this, we simply need to show that for all $z\in \Zc$, $f(\cdot,\cdot,z)$ is $\rm{PL}(2)$ (for some uniform PL constant). This can be shown as follows. Let $C=\sup_{z\in\Zc}|z|$. Then
\begin{align}
|f(x,y,z) - f(x',y',z)| &=|z(y-g(x))^2-z(y'-g(x'))^2|\\
%&\leq |z||(y-g(x))^2-(y'-g(x'))^2|\\
&\leq C|(y-g(x))^2-(y'-g(x'))^2|\\
&\leq C(1+\tn{\ub}+\tn{\vb})\tn{\ub-\vb},
 \end{align}
%\begin{align}
%|F_n(\btha_n(\g,\h),\betas,\Sigmab) - F_n(\bt,\betas,\Sigmab)| &=
% \frac{1}{p}\sum_{i=1}^p|\lai|\left| (\sqrt{p}\betas_i-g(\sqrt{p}\btha_{n,i}))^2-(\sqrt{p}\betas_i-g(\sqrt{p}\bt_{i}))^2 \right|\nn\\
% &\leq
%\Sigma_{\max} \frac{1}{p}\sum_{i=1}^p \left| (\sqrt{p}\betas_i-g(\sqrt{p}\btha_{n,i}))^2-(\sqrt{p}\betas_i-g(\sqrt{p}\bt_{i}))^2 \right|\nn\\
% &\leq
%{C} \Sigma_{\max} \frac{1}{p}\sum_{i=1}^p (1+ \|\sqrt{p}[\betas_i,\btha_{n,i}]\|_2 +  \|\sqrt{p}[\betas_i,\bt_{i}]\|_2) \sqrt{p}|\btha_{n,i}-\bt_{i}|\nn\\
% &\leq
%{C} \Sigma_{\max}  \Big(1+ \frac{1}{\sqrt{p}}\big(\sum_{i=1}^{p}\|\sqrt{p}[\betas_i,\btha_{n,i}]\|_2^2\big)^{1/2} + \frac{1}{\sqrt{p}}\big(\sum_{i=1}^p\|\sqrt{p}[\betas_i,\bt_{i}]\|_2^2\big)^{1/2}\Big) \|\btha_n-\bt\|_2\nn\\
% &\leq
%%C  \left(1+ \|\betas\|_2^2+ \|\btha_n\|_2^2 + \|\bt\|_2^2\right) \|\btha_n-\bt\|_2.
%C  \left(1+ \max\{\|\betas\|_2^2,\|\btha_n\|_2^2,\|\bt\|_2^2\}^{1/2} \right) \|\btha_n-\bt\|_2.\label{eq:fcase1}
% \end{align}
where $\ub=[x,y]$ and $\vb=[x',y']$. In the second line above, we used boundedness of $\Zc$. In the third line, we used the fact that the function $\psi(x,y) = (y-g(x))^2$ is $\rm{PL}(2)$ as it is a quadratic of a Lipschitz function.
}

This completes the proof of the lemma.
\end{proof}

{
\begin{lemma}[PL with Bounded Variables]\label{lem:PLbdd}%with bounded derivatives 
Let $f:\R^{d_1}\rightarrow \R$ be a $PL(k)$ function, $\Mc\subset \R^{d_2}$ be a compact set and $g$ be a continuously differentiable function over $\Mc$. Then $h(\x,\y)=f(\x)g(\y)$ is $PL(k+1)$ over $\R^{d_1}\times \Mc$.
\end{lemma}
\begin{proof} First observe that $g$ has continuous derivatives and is continuous over a compact set. Thus $g$ and its gradient are bounded and $g$ is Lipschitz over $\Mc$. Let $B=\sup_{\x\in \Mc}\max |g(x)|,\tn{\nabla g(\x)}$. To proceed, given pairs $(\x,\y)$ and $(\x',\y')$ over $\R\times \Mc$, we have that
\begin{align}
|h(\x,\y)-h(\x',y')|&\leq |h(\x,\y)-h(\x',\y)|+ |h(\x',\y)-h(\x',\y')|\\
&\leq |f(\x)-f(\x')||g(\y)|+ |f(\x')||g(\y)-g(\y')|\\
&\leq B|f(\x)-f(\x')|+ B\tn{\y-\y'}|f(\x')|\\
&\leq B(1+\tn{\x'}^{k-1}+\tn{\x}^{k-1})\tn{\x-\x'}+ B\tn{\y-\y'}(1+\tn{\x'}^k)\\
&\leq C (1+\tn{\z}^k+\tn{\z'}^{k})\tn{\z-\z'},
\end{align}
where $\z=[\x~\y]^T$ and $C$ an absolute constant. This shows the desired $\text{PL}(k+1)$ guarantee.
\end{proof}
The following lemma is in similar spirit to Lemma \ref{lem:PLbdd} and essentially follows from similar lines of arguments (i.e.~using Lipschitzness induced by boundedness).
}
\begin{lemma}\label{lem:hPL}
Let functions $f,g:\R^3\rightarrow\R$ such that $f\in\rm{PL}(k)$ and 
$$
g(x,y,z) := \frac{y^{-1/2}}{1+(\ksi_* y)^{-1}}\sqrt{\kappa}\tau_* z + (1-(1+\ksi_*y)^{-1})x.
$$
Here, $\ksi_*,\tau_*,\kappa$ are positive constants. 
%bounded, say $x_2\in[m,M]\subset(0,\infty)$.
Further define 
\begin{align}
h(x,y,z) := f\left(g(y,z,x),y,z\right),
\end{align}
and assume that $y$ take values on a fixed bounded compact set $\Mc{\subset\R^+}$.
Then, it holds that $h\in\rm{PL}(k+1)$.
% \fx{and $h(x,y,z)$ is continuously differentiable over $y\in \Mc$}
\end{lemma}

\begin{proof}
Since $f$ is $\rm{PL}(k)$, for some $L>0$, \eqref{PL func} holds. Fix $x,x'\in\R$, $\ab=[y,z]\in\R^{2}$ and $\ab'=[y',z']\in\R^{2}$. Let $\bb=[\g,\ab]=[\g,y,z]\in\R^3$ where $\g=g(y,z,x)\in\R$ and define accordingly $\bb'=[\g',\ab']$ and $\g'=g(y',z',x')$. We have that
\begin{align}
|h([x,\ab])-h([x',\ab'])|&= |f(\bb)-f(\bb')|\nn\\
&\leq L\left(1+\|\bb\|_2^{k-1}+\|\bb'\|_2^{k-1}\right)\|\bb-\bb'\|_2\nn\\
&\leq C\left(1+\|\ab\|_2^{k-1}+\|\ab'\|_2^{k-1}+|\g|^{k-1}+|\g'|^{k-1}\right)(\|\ab-\ab'\|_2+|\g-\g'|),\label{eq:hlip}
%~~~~~~~~&L\left(1+\|\g(\y,\z,\x)\|_2^{k-1}+\|g(\y',\z',\x')\|_2^{k-1}\right)\|g(\y,\z,\x)-g(\y',\z',\x')\|_2.
\end{align}
for some constant $C>0$. In the last inequality we have repeatedly used the inequality
$
\left(\sum_{i=1}^m \|\vb_i\|_2^2\right)^{\frac{d}{2}} \leq  C(m)\cdot\sum_{i=1}^m\|\vb_i\|_2^{d}.
$
Next, we need to bound the $\g$ term in terms of $(x,\ab)$. This is accomplished as follows
\begin{align}
|\g|^{k-1} &= \left|\frac{y^{-1/2}}{1+(\ksi y)^{-1}}\sqrt{\kappa}\tau_* z + (1-(1+\ksi_*y)^{-1})x\right|^{k-1}\nn\\
%\|\g\|_2^{k-1}&\leq C \frac{1}{p}\sum_{i=1}^p|g_i|^{k-1}\\
%&\lesssim \frac{1}{p}\sum_{i=1}^p\left|\frac{y_i^{-1/2}}{1+(\ksi y_i)^{-1}}\sqrt{\kappa}\tau_* z_i + (1-(1+\ksi_*y_i)^{-1})x_i\right|\\
&\leq C (|z|+|x|)^{k-1} \nn\\
&\leq C \left(|x|^{k-1}+|z|^{k-1}\right) \leq C (|x|^{k-1} + \|\ab\|_2^{k-1}).\label{eq:gb}
\end{align}
Here, the value of the constant $C>0$ may change from line to line. 
Secondly and similarly, we have the following perturbation bound on the $\g-\g'$. Recall that $y,y'\subset \Mc$ are bounded. Additionally, since $\Mc\subset\R^+$ and is compact, $\Mc$ is strictly bounded away from $0$. Let 
\[
g_1(y)=\sqrt{\kappa}\tau_*\frac{y^{-1/2}}{1+(\ksi_* y)^{-1}}\quad\text{and}\quad g_2(y)=1-(1+\ksi_*y)^{-1}.
\]
It can be seen that $g_1,g_2$ are continuously differentiable functions over $\Mc$. Thus $g_1,g_2$ are bounded and have bounded derivatives over $\Mc$. We will prove the following sequence of inequalities
%\CT{Not sure why this is red. Looks ok to me} 
% \so{Suppose the triples $(x_i,y_i,z_i)$ takes values in a fixed bounded compact set $\Mc$. We will prove the following sequence of inequalities} 
\begin{align}
%|\|\g\|_2-\|\g'\|_2|&\leq 
|{\g-\g'}|&= |g(y,z,x)-g(y',z',x')| \nn\\
&\leq \left|g_1(y) x - g_1(y') x'\right|
 + \left|g_2(y)z-g_2(y')z'\right|\nn\\
 &\leq
\left|g_1(y) x - g_1(y) x'\right| +\left|g_1(y) x'- g_1(y') x'\right|\nn\\
 &\qquad+ \left|g_2(y)z-g_2(y)z'\right| + \left|g_2(y)z'-g_2(y')z'\right| \nn\\
&\leq C_1|x-x'| + C_2|x'||y-y'| + C_3|z-z'| + C_4|z'||y-y'|\nn\\
&\leq C(1+|x'| + |z'|)(|x-x'| + |z-z'| + |y-y'|)\\
&\leq C\sqrt{3}(1+|x'| + |z'|)\|[\ab,x]-[\ab',x']\|_2
.\label{eq:gd}
%C^2_{\Mc} \sum_{i=1}^p(|x_i-x'_i|^2+|y_i-y'_i|^2+|z_i-z'_i|^2)\label{sec line eq}\\
%&\leq C^2_{\Mc} (\tn{\x-\x'}^2+\tn{\y-\y'}^2+\tn{\z-\z'}^2)\\
%&\lesssim \tn{\ab-\ab'}^2.
\end{align}
In the fourth inequality above, we used the fact that $|g_i(y)|,|g_i'(y)|$ are bounded. In the last line, we used Cauchy-Scwhartz.

Substituting \eqref{eq:gb} and \eqref{eq:gd} in \eqref{eq:hlip} gives:
\begin{align}
|h(x,y,z) - h(x',y',z')|  &\leq C\left(1+\|\ab\|_2^{k-1}+\|\ab'\|_2^{k-1}+|x|^{k-1}+|x'|^{k-1}\right)(1+|x'| + |z'|)\|[\ab,x]-[\ab',x']\|_2\nn \\
&\leq C\left(1+\|[\ab,x]\|_2^{k}+\|[\ab',x']\|_2^{k}\right)\|[\ab,x]-[\ab',x']\|_2.
\end{align}
Thus, $h\in\rm{PL}(k+1)$, as desired.
%In what follows, we prove the second line \eqref{sec line eq} i.e.~the fact that for any triples $a=(x,y,z),a'=(x',y',z')$ (with $a,a'\in\R^3$), we have that 
%\[
%|g(a)-g(a')|\leq C_{\Mc}\tn{a-a'}.
%\]
%Set $C_\nabla=\sup_{a\in \Mc}\tn{\nabla g(a)}$. By definition of gradient, the inequality above holds with $C_\Mc=C_\nabla$. Thus, all that remains is proving that $C_\nabla$ is upper bounded by a constant. \so{However, this automatically holds because from the definition of $g(\dots)$ function, it is clear that $C_\nabla$ is defined everywhere and continuous thus it has a finite maximum over a compact set.}
\end{proof}






